% ===== roadmap of this week =====
\begin{frame}{이번 주차 개요}
  \textbf{이 주차를 마치면}
  \begin{enumerate}
    \item 로봇공학 3대 기둥(상태공간 · 운동학 · 동역학)을 제 언어로 설명하고,
      ``시뮬레이션 한 스텝''이 미분방정식의 시간 적분임을 말할 수 있다.
    \item \kb{sim-to-real gap}가 왜 생기는지(노이즈가 아니라 시스템틱한 편차)
      설명하고, 도메인 무작위화와 시스템 ID가 서로 다른 시도임을 짚을 수 있다.
    \item 연속 제어 평가 프로토콜(창 고정, 여러 시드 평균)을 스스로 세우고,
      PD 제어 베이스라인과 정당한 비교를 할 수 있다.
    \item Week 11의 PPO를 이산 행동에서 \kb{연속(가우시안) 정책}으로
      확장해 Reacher 환경에서 목표 추적 정책을 학습시킬 수 있다.
  \end{enumerate}
  \vskip0.6em
  \small
  \begin{itemize}
    \item \kb{13.1} 2관절 팔 FK/IK 손계산 · 진자 동역학 · PD 제어 · 현실 격차
    \item \kb{13.2} PD 그리드 탐색 · 평가 프로토콜 · Reacher/Ant · 4항 장부
    \item \kb{13.3} 가우시안 정책 · Reacher 10차원 상태의 내면 · PPO 100만 스텝
  \end{itemize}
\end{frame}
